% !TEX TS-program = pdflatex
% !TeX program = pdflatex
% !TEX encoding = UTF-8
% !TEX spellcheck = en

\documentclass[11pt, a4paper]{article}
%\usepackage{fullpage}
\usepackage[left=1.5cm,right=1.5cm,top=1cm,bottom=2cm]{geometry}
\usepackage[fleqn]{amsmath}
\usepackage{amssymb}
%\usepackage{indentfirst}
\usepackage[T1]{fontenc}
\usepackage[utf8]{inputenc}
\usepackage[french,english]{babel}
\usepackage{txfonts}
\usepackage[]{graphicx}
\usepackage{multirow}
\usepackage{hyperref}
\usepackage{parskip}
\usepackage{multicol}
\usepackage{wrapfig}
\usepackage{fancyhdr}
\usepackage{hyperref}
\usepackage{tcolorbox}
\usepackage{natbib}

\usepackage{turnstile}%Induction symbole

\renewcommand{\baselinestretch}{1}

\setlength{\parindent}{24pt}

%\usepackage{etoolbox}
%\patchcmd{\thebibliography}{\section*{\refname}}{}{}{}

\fancyhf{}

\lfoot{Dr. Abdelkrime Aries}
\cfoot{2024-2025/ESI/Master/NLP}
\rfoot{\textbf{\thepage}}
%
\renewcommand{\headrulewidth}{0pt} % remove lines as well
\renewcommand{\footrulewidth}{0pt}

\newcommand\repeatstr[1]{\leavevmode\xleaders\hbox{#1}\hfill\kern0pt}

\renewcommand{\bibsection}{}
\bibliographystyle{humannat}

\newcommand{\kurl}[1]{{\scriptsize\bfseries\url{#1}}}

\setlength{\parindent}{0pt}

\begin{document}

%\selectlanguage {french}
%\pagestyle{empty} 
\pagestyle{fancy}

\noindent
\begin{tabular}{ll}
\multirow{3}{*}{\includegraphics[width=1.6cm]{../../extra/logo/esi.nlp.pdf}} & \'Ecole national Supérieure d'Informatique, Algiers\\
& Master 2024-2025\\
& Natural Language Processing (NLP)
\end{tabular}\\[.25cm]
\noindent\rule{\textwidth}{1pt}\\[-0.25cm]
\begin{center}
{\LARGE \textbf{Syllabus}}
\begin{flushright}
	Dr. Abdelkrime Aries
\end{flushright}
\end{center}\vspace*{-0.25cm}
\noindent\rule{\textwidth}{1pt}

\begin{abstract}
	Natural Language Processing (NLP) is an interdisciplinary field involving linguistics, computer science, and artificial intelligence. 
	It is a domain of research and computer application dedicated to the understanding and generation of human language by computers. 
	With the emergence of language models such as OpenAI's GPT, Google's BERT, and others, NLP has become a rapidly growing field with increasing importance in various industries and domains.
	
	NLP covers a wide range of topics, and it would be impossible to cover them all in a week. 
	However, in this seminar, we will focus on the most relevant and current tasks in scientific research.
\end{abstract}

\section{General information}

\begin{minipage}{0.44\textwidth}
		
	\subsection{Prerequisites}
	
	\begin{itemize}
		\item Basic knowledge of languages (vocabulary, grammar, etc.);
		\item Data structures (1CP-ALSDD);
		\item Probability and statistics \\(2CP-PRST1, 2CP-PRST2);
		\item Algebra and calculus \\(2CP-ANAL3, 2CP-ALG3, 2CP-ANAL4);
		\item Language theory (1CS-THP);
		\item Operations research (1CS-RO);
		\item Machine learning (2CS-ML or 2CS-BDM).
	\end{itemize}

\end{minipage}
\begin{minipage}{0.05\textwidth}
	\hspace{2pt}
\end{minipage}
\begin{minipage}{0.49\textwidth}
	\subsection{Objectives}
	
	\begin{itemize}
		\item Learn some linguistic concepts;
		\item Apply concepts covered in the prerequisite modules to solve NLP tasks;
		\item Discover some NLP tools and resources;
		\item Understand the research directions in this field;
		\item Acquire the ability to articulate a novel NLP methodology and assess its viability within the framework of a scholarly publication.
	\end{itemize}
	
\end{minipage}

\section{Outline}

\subsection{Lecture}

\begin{enumerate}
	\item \textbf{Introduction}: History, Language processing levels, Applications of NLP, NLP Challenges.
	\item \textbf{Machine learning for NLP}: Machine learning (review), Text classification, Sequences classification, Attention.
	\item \textbf{Basic text processing}: Characters, Text segmentation, Text normalization and filtering, Morphology.
	\item \textbf{Language models}: N-gram model, Neural models, Evaluation.
	\item \textbf{Part-of-speech tagging}: Sequences labeling, Task description, Approaches.
	\item \textbf{Word semantics}: Lexical databases, Vector representation of words, Word embedding, Word Sense Disambiguation (WSD).
	\item \textbf{Some applications}: discussion of some research works in this field.
\end{enumerate}

\subsection{Practical}

\begin{itemize}
	\item \textbf{Lab01}: Discovery of some NLP tools.
	The purpose of this lab is to assess the student's ability to follow instructions in order to create an application. Additionally, students will learn how to use some NLP tools.
	A perfect solution is one that precisely follows the lab's instructions.
	\item \textbf{Lab02}: Developing an NLP application.
	The purpose of this practical work is to test the student's ability to enhance an existing application based on some metrics.
	\item \textbf{Lab03}: Writing an article on this application.
	The purpose of this practical work is to teach students how to write an article in the field of NLP and showcase their work. 
	Additionally, students will learn how to use \LaTeX.
\end{itemize}
%\clearpage

The practical work can be done individually or in collaboration with a partner. Before participating in the practical work, it is essential to have a personal computer and a working environment, either by using online Colab (requiring an internet connection) or by having Anaconda installed.
In case of choosing Anaconda, you will need to install other tools like NLTK and TensorFlow. A \LaTeX\ editing environment is also required, either by using Overleaf (which requires an internet connection) or by installing \LaTeX\footnote{\LaTeX: \kurl{https://www.latex-project.org/get/}} on your machine.

\section{Evaluation}

Students will be evaluated using three methods:
\begin{itemize}
	\item \textbf{Quizzes (20\%)}: There will be at least 3 online quizzes. 
	Each quiz will cover the chapters seen during the same day.
	It must be completed individually before the end of the day (23h59).
	\item \textbf{Labs (30\%)}: The labs are already described in "Practical" section.
	\item \textbf{Written Exam (50\%)}: The exam is the most important part of evaluation. 
	Being absent on the day of the exam means a disqualification from this course.
\end{itemize}

\section{Bibliography}

\subsection{Similar courses}

\begin{itemize}
	\item "\textit{NLP S21}" (up to date).
	Carnegie Mellon Univ. 
	\kurl{http://demo.clab.cs.cmu.edu/NLP/}
	
	\item "\textit{CS388: NLP}" (2018). 
	Univ. of Texas. 
	\kurl{https://www.cs.utexas.edu/~mooney/cs388/}
	
	\item "\textit{CSE 517: NLP}" (NA).
	Univ. of Washington.
	\kurl{https://courses.cs.washington.edu/courses/cse517/}
	
	\item "\textit{CS 294-5: Statistical NLP}" (2005). 
	Univ. of Berkely. 
	\kurl{https://people.eecs.berkeley.edu/~klein/cs294-5/index.html}
	
	\item "\textit{CS224n: NLP with Deep Learning}" (up to date).
	Stanford Univ.
	\kurl{https://web.stanford.edu/class/cs224n/}
	
	\item "\textit{CS 5340/6340: NLP}" (up to date). 
	Univ. of Utah.
	\kurl{https://my.eng.utah.edu/~cs5340/schedule.html}
\end{itemize}

\subsection{To read (NLP)}

\begin{itemize}
	\item "\textit{TALN - course (ESI)}" [FR].
	Abdelkrime Aries. 
	\kurl{https://github.com/projeduc/ESI_2CS_TALN/releases/download/4.0/taln_v2.pdf}
	
	
	\item "\textit{Speech and Language Processing}".
	Dan Jurafsky and James H. Martin.
	\kurl{https://web.stanford.edu/~jurafsky/slp3/}
	
	\item "\textit{NLP}".
	Jacob Esenstein.
	\kurl{https://cseweb.ucsd.edu/~nnakashole/teaching/eisenstein-nov18.pdf}
	
	\item "\textit{Modern Approaches in NLP - course (LMU Munich)}".
	Carolin Becker, et al.
	\kurl{https://slds-lmu.github.io/seminar_nlp_ss20/index.html}
\end{itemize}

\subsection{To read (ML)}

\begin{itemize}
	\item "\textit{Dive into Deep Learning}" [online].
	Many authors (Amazon).
	\kurl{https://d2l.ai/index.html}
	
	\item "\textit{Deep Learning - course (NYU Center for data science)}".
	Yann LeCun and Alfredo Canziani.
	\kurl{https://atcold.github.io/pytorch-Deep-Learning/}
\end{itemize}

\end{document}
